\documentclass[11pt]{article}
\usepackage[a4paper,margin=21mm]{geometry}
\usepackage{fontspec}
\setmainfont{Latin Modern Roman}
\newfontfamily\CJKfont{Noto Serif CJK SC}
\XeTeXlinebreaklocale "zh"
\XeTeXlinebreakskip=0pt plus 1pt
\usepackage{amsmath,amssymb,amsthm,amscd,graphicx,color}
\usepackage{xurl}
\usepackage[unicode,hidelinks]{hyperref}
\allowdisplaybreaks
\setlength\emergencystretch{3em}
\numberwithin{equation}{section}


\numberwithin{equation}{section}

\newtheorem{Theorem}{Theorem}[section]
\newtheorem*{Theorem*}{Theorem}
\newtheorem{Corollary}[Theorem]{Corollary}
\newtheorem{Lemma}[Theorem]{Lemma}
\newtheorem{Proposition}[Theorem]{Proposition}
 { \theoremstyle{definition}
\newtheorem{Definition}[Theorem]{Definition}
\newtheorem{Note}[Theorem]{Note}
\newtheorem{Example}[Theorem]{Example}
\newtheorem{Remark}[Theorem]{Remark}
\newtheorem{Remarks}[Theorem]{Remarks}
}

\begin{document}
\begin{center}\Large Euclidean boundary-immersion integral inequality and conditional flat-ball rigidity under nonnegative p-form curvature and positive (n+1-p)-curvature\end{center}
\noindent\textbf{Stable ID:} 1169\quad\textbf{Author:} Nicolas Ginoux; Georges Habib; Simon Raulot\par
\noindent\textbf{Work:} A Poincar\'e Formula for Differential Forms and Applications\par
\noindent\textbf{Source:} \url{https://sigma-journal.com/2023/088/}\par
\noindent\textbf{Version:} Official SIGMA 19 (2023) article 088, DOI 10.3842/SIGMA.2023.088; journal PDF and native TeX acquired 2026-09-30, exact bytes pinned by SHA256\par
\noindent\textbf{Rights:} CC BY 4.0. Authors retain copyright. \url{https://creativecommons.org/licenses/by/4.0/}. Reuse and typesetting adaptation; original language and mathematical content preserved.\par
\noindent\textbf{Difficulty (prior selection preserved):} {\CJKfont H2: The author combines harmonic extensions of exact boundary p-forms with the form-valued Reilly identity and a complete exterior-power trace computation. The equality proof then constructs a maximal family of parallel p-forms through a nontrivial rank argument using the immersion second fundamental form. This forces ambient flatness and boundary umbilicity under the stated codimension-one and connectivity restrictions, giving the full ball rigidity rather than only a numerical inequality.}\par
\medskip\noindent\textbf{Editorial boundary note (not author text):} Complete original Theorem3.3 including every Euclidean immersion, p-form curvature, positive boundary-curvature and conditional equality/rigidity hypothesis. Full preceding Theorem3.1 harmonic-extension/Reilly-square argument and Remark3.2 are retained, followed by the complete selected sum/trace computation, boundary constancy and full rank/parallel-form proof. Exact named harmonic boundary-value/cohomology, Savo immersion identities and classical flat-ball/extension rigidity results remain precisely stated author references. The displayed first orthogonality computation and its explicitly parallel second calculation are retained without an editor derivation. One independently announced Euclidean boundary outcome; supporting form inequality and later spherical/Ros variants receive no extra counts. All literal signs, indices, conventions and order are preserved.\par
\medskip\hrule\medskip\noindent\textbf{Author text begins below.}\par
\setcounter{section}{1}
% Original source lines 175--411
\section{Preliminaries and notations}\label{s:prelim}

In this section, we briefly introduce the geometric setting and fix the notations of this paper.

Let $\big(M^{n+1},g\big)$ be a compact oriented $(n+1)$- dimensional Riemannian ma\-ni\-fold with smooth nonempty boundary $\partial M$ and let $\iota\colon\partial M\to M$ be the inclusion map.
Let ${\rm d}\mu_g$ be the Riemannian measure induced by $g$ on both $M$ and $\partial M$.
In the following, any metric -- being $g$ on $TM$ or any metric induced by $g$ on further bundles -- will be denoted by $\langle\cdot\,,\cdot\rangle$, with associated pointwise norm $|\cdot|$.
Let $\nu$ denote the inward unit normal along $\partial M$ and $S:=-\nabla\nu$ be the associated Weingarten endomorphism-field, where $\nabla$ denotes the Levi-Civita connection on $TM$.
Let $H:=\frac{1}{n}\operatorname{tr}(S)$ denote the mean curvature of $\partial M$ in $M$.
For any integer $p\in\{0,\dots,n+1\}$, let $\Omega^p(M):=\Gamma(\Lambda^pT^*M)$ be the space of differential forms of degree $p$ on $M$, that is the space of sections of the exterior bundle $\Lambda^p T^*M\to M$.
Let $\star\colon\Lambda^pT^*M\to\Lambda^{n+1-p}T^*M$ denote the pointwise Hodge star  operator.
For any $(1,1)$-tensor $A$ and $p$-form $\omega$, let $A^{[p]}\omega$ be the $p$-form that is pointwise defined by the following identity: for all tangent vectors $X_1,\dots,X_p$,
\[\bigl(A^{[p]}\omega\bigr)(X_1,\dots,X_p):=\sum_{j=1}^p\omega(X_1,\dots,AX_j,\dots,X_p).\]
In the particular case, where $A=S$, we denote for each $k\in\{1,\dots,n\}$ by $\sigma_k$ the pointwise $k$-curvature of $\partial M$, that is the sum of the $k$ smallest principal curvatures (i.e., the eigenvalues of~$S$) of $\partial M$.
Note that, since the eigenvalues of $S$ are the sums of exactly $p$ among the $n$~principal curvatures,
\[\bigl\langle S^{[p]}\omega,\omega\bigr\rangle\geq\sigma_p|\omega|^2\]
for any $\omega\in\Lambda^pT^*\partial M$, with equality if and only if $S^{[p]}\omega=\sigma_p\omega$.
Moreover, for all $1\leq p\leq q\leq n$, we have $\frac{\sigma_p}{p}\leq\frac{\sigma_q}{q}$, with equality when $p<q$ if and only if the $q$ smallest principal curvatures are equal.

Let ${\rm d}$ (resp.\ $\delta$) denote the exterior derivative (resp.\ codifferential) on $p$-forms and $\nabla$ be the covariant derivative induced by $\nabla$ on $\Lambda^pT^*M$.
Recall the so-called \emph{Reilly formula} \cite[Theorem~3]{RaulotSavo11} for differential $p$-forms with $p\geq 1$: for any $\omega\in\Omega^p(M)$
\begin{gather}
\int_M\bigl(|{\rm d}\omega|^2+|\delta\omega|^2-|\nabla\omega|^2-\bigl\langle W^{[p]}\omega,\omega\bigr\rangle\bigr)\,{\rm d}\mu_g \nonumber\\
\qquad{} =\int_{\partial M}\bigl(2\bigl\langle\nu\lrcorner\,\omega,\delta^{\partial M}(\iota^*\omega)\bigr\rangle+\bigl\langle S^{[p]}\iota^*\omega,\iota^*\omega\bigr\rangle+\bigl\langle S^{[n+1-p]}\iota^*(\star\omega),\iota^*(\star\omega)\bigr\rangle\bigr)\,{\rm d}\mu_g,\label{eq:Reillyformula}
\end{gather}
where $\delta^{\partial M}$ denotes the codifferential on $\partial M$ and $W^{[p]}$ the curvature term involved in the Weitzenb\"ock formula for $p$ forms: denoting by $\Delta:={\rm d}\delta+\delta {\rm d}$ the Hodge Laplace operator on $p$-forms,
\[\Delta\omega=\nabla^*\nabla\omega+W^{[p]}\omega\]
for any $p$-form $\omega$.  By convention, we let $W^{[p]}=0$ and $S^{[p]}=0$ for all $p\leq 0$.
Note that \cite[Theorem 3]{RaulotSavo11}
\begin{equation*}%\label{Sstaromega}
\bigl\langle S^{[n+1-p]}\iota^*(\star\omega),\iota^*(\star\omega)\bigr\rangle=nH|\nu\lrcorner\,\omega|^2-\bigl\langle S^{[p-1]}(\nu\lrcorner\,\omega),\nu\lrcorner\,\omega\bigr\rangle
\end{equation*}
for any $p$-form $\omega$.

\section[A Poincar\'e-type inequality for p-forms]{A Poincar\'e-type inequality for $\boldsymbol p$-forms}\label{s:Poincareineqpformsgeneral}

We first prove a generalized version of the integral i\-ne\-qua\-li\-ty for functions obtained by Miao and Wang in \cite[equation~(1.3)]{MiaoWang14}.

\begin{Theorem}\label{t:Poincareineqpformsgeneral}
Let $\bigl(M^{n+1},g\bigr)$ be any compact oriented  $(n+1)$-dimensional Riemannian manifold with smooth nonempty boundary $\partial M$.
Fix $p\in\{1,\dots,n\}$ and assume that $W^{[p]}\geq0$ on $M$ as well as $\sigma_{n+1-p}>0$ along $\partial M$.
Then, for any exact $p$-form $\omega$ on $\partial M$,  we have
\begin{equation}\label{eq:Poincareineqpformsgeneral}
\int_{\partial M}\frac{\bigl|\delta^{\partial M}\omega\bigr|^2}{\sigma_{n+1-p}}\,{\rm d}\mu_g\geq\int_{\partial M}\langle S^{[p]}\omega,\omega\rangle\,{\rm d}\mu_g.
\end{equation}
Moreover, if {\rm (\ref{eq:Poincareineqpformsgeneral})} is an equality for some non-zero $\omega$, then for the $p$-form $\hat{\omega}$ on $M$ satisfying ${\rm d}\hat{\omega}=0=\delta\hat{\omega}$ on $M$ as well as $\iota^*\hat{\omega}=\omega$ on $\partial M$, the identities $\delta^{\partial M}\omega=-\sigma_{n+1-p}\nu\lrcorner\,\hat{\omega}$, $S^{[p-1]}(\nu\lrcorner\,\hat{\omega})=(nH-\sigma_{n+1-p})\nu\lrcorner\,\hat{\omega}$ hold along $\partial M$ and the $p$-form $\hat{\omega}$ must be parallel -- hence $W^{[p]}\hat{\omega}=0$ must hold -- on $M$.
\end{Theorem}

\begin{proof}
The proof mainly follows that of \cite[Theorem 5]{RaulotSavo11}.
Let $\omega$ be an exact $p$-form on $\partial M$, that is $\omega={\rm d}^{\partial M}\alpha$ for some $\alpha\in\Omega^{p-1}(\partial M)$.
By \cite[Theorem 2]{DuffSpencer52}, there exists a $(p-1)$-form $\hat{\alpha}$ on~$M$ such that $\delta {\rm d}\hat{\alpha}=0$ on $M$ with $\iota^*\hat{\alpha}=\alpha$ along $\partial M$. Let $\hat{\omega}:={\rm d}\hat{\alpha}\in\Omega^p(M)$, then $\hat{\omega}$ satisfies ${\rm d}\hat{\omega}=0$, $\delta\hat{\omega}=\delta {\rm d}\hat{\alpha}=0$ on~$M$ with $\iota^*\hat{\omega}=\iota^*({\rm d}\hat{\alpha})=\omega$ along $\partial M$. Actually such a~$p$-form~$\hat{\omega}$ on~$M$ with ${\rm d}\hat{\omega}=0=\delta\hat{\omega}$ on $M$ as well as $\iota^*\hat{\omega}=\omega$ along $\partial M$ is uniquely determined by $\omega$. Namely, because of $W^{[p]}\geq0$ by assumption and the property $\star W^{[p]}\star^{-1}=W^{[n+1-p]}$, we know that $W^{[n+1-p]}\geq0$. Together with the assumption $\sigma_{n+1-p}>0$ we can deduce from \cite[Theorem~4]{RaulotSavo11} that $H_{\rm abs}^{n+1-p}(M)=0$  holds, where
\[H_{\rm abs}^{k}(M) :=\bigl\{\alpha\in\Omega^k(M) \mid {\rm d}\alpha=0=\delta\alpha\textrm{ and }\nu\lrcorner\,\alpha=0\bigr\}\]
is, for every $0\leq k\leq n+1$, the $k^{\textrm{th}}$ absolute de Rham cohomology group. By Poincar\'e duality, $H_{\rm abs}^{k}(M)\cong H_{\rm rel}^{n+1-k}(M)$,  where
\[H_{\rm rel}^{k}(M):=\bigl\{\alpha\in\Omega^k(M) \mid {\rm d}\alpha=0=\delta\alpha\textrm{ and }\iota^*\alpha=0\bigr\}\]
is the $k^{\textrm{th}}$ relative de Rham cohomology group. Therefore, $H_{\rm rel}^{p}(M)=0$, so that $\hat{\omega}$ is uniquely determined by $\omega$.

We apply identity (\ref{eq:Reillyformula}) to $\hat{\omega}$: since ${\rm d}\hat{\omega}=\delta\hat{\omega}=0$ by construction of $\hat{\omega}$ and $\bigl\langle W^{[p]}\hat{\omega},\hat{\omega}\bigr\rangle\geq0$ by assumption,  as well as $|\nabla\hat{\omega}|^2\geq 0$, we have from (\ref{eq:Reillyformula}) that
\[0\geq\int_{\partial M}\bigl(2\bigl\langle\nu\lrcorner\,\hat{\omega},\delta^{\partial M}\omega\bigr\rangle+\bigl\langle S^{[p]}\omega,\omega\bigr\rangle+\bigl\langle S^{[n+1-p]}\iota^*(\star\hat{\omega}),\iota^*(\star\hat{\omega})\bigr\rangle\bigr)\,{\rm d}\mu_g.\]
By definition of the $k$-curvatures and the identity $\iota^*(\star\hat{\omega})=(-1)^p\nu\lrcorner\star(\nu\lrcorner\,\hat{\omega})=\star_{\partial M}(\nu\lrcorner\,\hat{\omega})$, we have\looseness=-1
\begin{equation}\label{eq:sn1p}
\bigl\langle S^{[n+1-p]}\iota^*(\star\hat{\omega}),\iota^*(\star\hat{\omega})\bigr\rangle\geq\sigma_{n+1-p}|\iota^*(\star\hat{\omega})|^2=\sigma_{n+1-p}|\nu\lrcorner\,\hat{\omega}|^2
\end{equation}
on $\partial M$. Moreover, because $\sigma_{n+1-p}$ is assumed to be positive along $\partial M$,  we write
\begin{gather*}
2\bigl\langle\nu\lrcorner\,\hat{\omega},\delta^{\partial M}\omega\bigr\rangle=\Biggl|\sigma_{n+1-p}^{1/2}\nu\lrcorner\,\hat{\omega}+\frac{1}{\sigma_{n+1-p}^{1/2}}\delta^{\partial M}\omega\Biggr|^2-\sigma_{n+1-p}|\nu\lrcorner\,\hat{\omega}|^2-\frac{|\delta^{\partial M}\omega|^2}{\sigma_{n+1-p}}\\
\hphantom{2\bigl\langle\nu\lrcorner\,\hat{\omega},\delta^{\partial M}\omega\bigr\rangle}
\geq-\sigma_{n+1-p}|\nu\lrcorner\,\hat{\omega}|^2-\frac{|\delta^{\partial M}\omega|^2}{\sigma_{n+1-p}},
\end{gather*}
so that
\[0\geq \int_{\partial M}\left(-\frac{|\delta^{\partial M}\omega|^2}{\sigma_{n+1-p}}+\bigl\langle S^{[p]}\omega,\omega\bigr\rangle\right)\,{\rm d}\mu_g\]
which is inequality   (\ref{eq:Poincareineqpformsgeneral}).

Assume now (\ref{eq:Poincareineqpformsgeneral}) to be an equality for some non-zero $\omega$. Let $\hat{\omega}$ be the $p$-form on $M$ such that ${\rm d}\hat{\omega}=0=\delta\hat{\omega}$ on $M$ and $\iota^*\hat{\omega}=\omega$ on $\partial M$; as mentioned above.
By the sequence of pointwise i\-ne\-qua\-li\-ties used in the above proof of i\-ne\-qua\-li\-ty (\ref{eq:Poincareineqpformsgeneral}), we can deduce that $\nabla\hat{\omega}=0$ on $M$, that is,~$\hat{\omega}$ is parallel on $M$  and, furthermore, $\delta^{\partial M}\omega=-\sigma_{n+1-p}\nu\lrcorner\,\hat{\omega}$ and $S^{[n+1-p]}\iota^*(\star\hat{\omega})=\sigma_{n+1-p}\iota^*(\star\hat{\omega})$ must hold on $\partial M$.  As a straightforward consequence, $W^{[p]}\hat{\omega}=0$ on $M$. Thanks to the identity $S^{[n+1-p]}(\star_{\partial M}\alpha)=-\star_{\partial M}S^{[p-1]}\alpha+nH\star_{\partial M}\alpha$ which is valid pointwise for all $(p-1)$-forms $\alpha$, the latter is equivalent to $S^{[p-1]}(\nu\lrcorner\,\hat{\omega})=(nH-\sigma_{n+1-p})\nu\lrcorner\,\hat{\omega}$ along $\partial M$. This concludes the proof of~Theorem \ref{t:Poincareineqpformsgeneral}.
\end{proof}


\begin{Remark}\label{r:LambdapTMtrivial} Note that, for $1\leq p\leq n$, if the bundle $\Lambda^p T^*M\to M$ is trivialized by such parallel $p$-forms $\hat{\omega}$, then on the one hand the manifold $M$ is flat and, on the other hand, $S^{[p-1]}=(nH-\sigma_{n+1-p})\cdot\mathrm{Id}$ must hold pointwise on $\Lambda^{p-1}T^*\partial M$.
The first statement is a consequence from the fact that the curvature of the manifold $M$ vanishes as soon as that of $\Lambda^p T^*M$ does.
The second statement comes from the map $\Lambda^p T_x^*M\to \Lambda^p T^*_x\partial M\oplus \Lambda^{p-1} T^*_x\partial M$, $\omega\mapsto (\iota^*\omega,\nu\wedge(\nu\lrcorner\,\omega))$, being an isomorphism at any $x\in \partial M$.  In case $p\geq2$ (for $p=1$ that identity is trivial because of $S^{[0]}=0$ by convention and $\sigma_{n}=nH$ by definition), this shows that $\iota\colon\partial M\to M$ must be \emph{totally umbilical}.
\end{Remark}

Next we turn to the case where the boundary of $M$ is assumed to be isometrically immersed in some Euclidean space.

\begin{Theorem}\label{t:PoincareineqpformsboundaryinEuclideanspace}
Let $\bigl(M^{n+1},g\bigr)$ be any compact oriented $(n+1)$-dimensional Riemannian manifold with smooth nonempty boundary $\partial M$.
Assume that $W^{[p]}\geq0$ on $M$ as well as $\sigma_{n+1-p}>0$ along $\partial M$ for a given $p\in\{1,\dots,n\}$.
Assume also that there exists an isometric immersion $\iota_0\colon\partial M\to\mathbb{R}^{n+m}$ with mean curvature vector $H_0$. Then
\begin{equation}\label{eq:PoincareineqpformsboundaryinEuclideanspace}
\int_{\partial M}\frac{n|H_0|^2-\frac{p-1}{n(n-1)}{\rm Scal}^{\partial M}}{\sigma_{n+1-p}}\,{\rm d}\mu_g\geq\int_{\partial M}H\,{\rm d}\mu_g,
\end{equation}
where  ${\rm Scal}^{\partial M}$  denotes the scalar curvature of $\partial M$. Equality holds in {\rm (\ref{eq:PoincareineqpformsboundaryinEuclideanspace})} when $M$ is the $(n+1)$-dimensional flat disk $\mathbb{D}^{n+1}$ standardly embedded in some $(n+1)$-dimensional affine subspace of~$\mathbb{R}^{n+m}$. Conversely, if {\rm (\ref{eq:PoincareineqpformsboundaryinEuclideanspace})} is an equality, $\partial M$ is connected, $p\geq 2$ and $\iota_0(\partial M)$ is contained in some $(n+1)$-dimensional affine subspace of $\mathbb{R}^{n+m}$, then, up to rescaling the metrics on $M$ and on $\mathbb{R}^{n+m}$, the manifold $M$ is isometric to the $(n+1)$-dimensional flat disk $\mathbb{D}^{n+1}$ standardly embedded in that subspace.
\end{Theorem}

\begin{proof}
We take the standard coordinates $(x_1,\dots,x_{n+m})$ on $\mathbb{R}^{n+m}$ and, for any $i_1,\dots,i_p\in\{1,\dots,n+m\}$, we denote by ${\rm d}x_I:={\rm d}x_{i_1}\wedge\dots\wedge {\rm d}x_{i_p}\in\Lambda^p(\mathbb{R}^{n+m})^*$ and by $\omega_I:=\iota_0^* {\rm d}x_I\in\Omega^p(\partial M)$. Note that, since ${\rm d}x_I$ is exact, so is $\omega_I$.
Replacing $\omega$ by $\omega_I$ in (\ref{eq:Poincareineqpformsgeneral}), we obtain
\begin{equation}\label{eq:PoincareomegaI}
\int_{\partial M}\frac{\bigl|\delta^{\partial M}\omega_I\bigr|^2}{\sigma_{n+1-p}}\,{\rm d}\mu_g\geq\int_{\partial M}\bigl\langle S^{[p]}\omega_I,\omega_I\bigr\rangle\,{\rm d}\mu_g.
\end{equation}
We now want to deduce from (\ref{eq:PoincareomegaI}) a more explicit i\-ne\-qua\-li\-ty. For this, we sum (\ref{eq:PoincareomegaI}) over $I$, meaning that we compute the sum $\sum_{i_1,\dots,i_p=1}^{n+m}$ of both sides; mind that the indices $i_1,\dots,i_p$ vary independently and hence are repeated. On the one hand, by \cite[Lemma 2.2]{Savo05},
\begin{equation}\label{eq:sumdeltaomegaI2}
\sum_I\bigl|\delta^{\partial M}\omega_I\bigr|^2=p!\binom{n}{p}p\left(n|H_0|^2-\frac{p-1}{n(n-1)}{\rm Scal}^{\partial M}\right).
\end{equation}
On the other hand, we use the pointwise identity
\[\sum_{j}{\rm e}_j\wedge({\rm e}_j\lrcorner\,\alpha)=k\alpha\]
which is valid for any $k$-form $\alpha$ and any pointwise o.n.b.\ $({\rm e}_j)_j$ of $TM$ or $T\partial M$.
As a straightforward consequence,
\begin{equation*}
\sum_{j=1}^n\langle {\rm e}_j\wedge\alpha,{\rm e}_j\wedge\beta\rangle=(n-k)\langle\alpha,\beta\rangle
\end{equation*}
for any $k$-forms $\alpha$, $\beta$ on an $n$-dimensional space. Thus we may compute the sum of the r.h.s.\ of~(\ref{eq:PoincareomegaI}) as follows:  Let $A\colon T\partial M\to T\partial M$ be any symmetric endomorphism of $\partial M$  and denoting by ${\rm d}x_i^T:=\iota_0^*{\rm d}x_i$ and by $({\rm e}_j)_{1\leq j\leq n}$ a pointwise o.n.b.\ of $T\partial M$, we may write
\begin{align}
\sum_I\bigl\langle A^{[p]}\omega_I,\omega_I\bigr\rangle&= \sum_{i_1,\dots,i_p=1}^{n+m}\bigl\langle A^{[p]}({\rm d}x_{i_1}^T\wedge\dots\wedge {\rm d}x_{i_p}^T),{\rm d}x_{i_1}^T\wedge\dots\wedge {\rm d}x_{i_p}^T\bigr\rangle\nonumber\\
&= \sum_{i_1=1}^{n+m}\sum_{i_2,\dots,i_p=1}^{n+m}\sum_{j,k=1}^n\bigl\langle {\rm d}x_{i_1}^T,{\rm e}_j\bigr\rangle\bigl\langle {\rm d}x_{i_1}^T,{\rm e}_k\bigr\rangle\bigl\langle A^{[p]}({\rm e}_j\wedge {\rm d}x_{i_2,\dots,i_p}^T),{\rm e}_k\wedge {\rm d}x_{i_2,\dots,i_p}^T\bigr\rangle\nonumber\\
&= \sum_{j,k=1}^n\!\underbrace{\left(\sum_{i_1=1}^{n+m}\!\bigl\langle {\rm d}x_{i_1},{\rm e}_j\bigr\rangle\bigl\langle {\rm d}x_{i_1},{\rm e}_k\bigr\rangle\!\right)\!}_{\delta_{jk}}\sum_{i_2,\dots,i_p=1}^{n+m}\!\!\bigl\langle A^{[p]}({\rm e}_j\wedge {\rm d}x_{i_2,\dots,i_p}^T),{\rm e}_k\wedge {\rm d}x_{i_2,\dots,i_p}^T\bigr\rangle\nonumber\\
&= \sum_{j=1}^n\sum_{i_2,\dots,i_p=1}^{n+m}\bigl\langle A^{[p]}({\rm e}_j\wedge {\rm d}x_{i_2,\dots,i_p}^T),{\rm e}_j\wedge {\rm d}x_{i_2,\dots,i_p}^T\bigr\rangle\nonumber\\
&= \sum_{j_1,\dots,j_p=1}^n\bigl\langle A^{[p]}({\rm e}_{j_1}\wedge\dots\wedge {\rm e}_{j_p}),{\rm e}_{j_1}\wedge\dots\wedge {\rm e}_{j_p}\bigr\rangle\nonumber\\
&= p\sum_{j_1,\dots,j_p=1}^n\bigl\langle A{\rm e}_{j_1}\wedge {\rm e}_{j_2}\wedge\dots\wedge {\rm e}_{j_p},{\rm e}_{j_1}\wedge\dots\wedge {\rm e}_{j_p}\bigr\rangle\nonumber\\
&= p(n-p+1)\sum_{j_1,\dots,j_{p-1}=1}^n\bigl\langle A{\rm e}_{j_1}\wedge {\rm e}_{j_2}\wedge\dots\wedge {\rm e}_{j_{p-1}},{\rm e}_{j_1}\wedge\dots\wedge {\rm e}_{j_{p-1}}\bigr\rangle\nonumber\\
&= p(n-p+1)\cdots(n-1)\sum_{j_1=1}^n\bigl\langle A{\rm e}_{j_1},{\rm e}_{j_1}\bigr\rangle\nonumber\\
&= p!\binom{n}{p}\frac{p}{n}{\rm tr}(A).\label{eq:traceap}
\end{align}
When $A=S$ is the associated Weingarten endomorphism-field, we get that
\begin{equation}\label{eq:sumISpomegaI}
\sum_I\bigl\langle S^{[p]}\omega_I,\omega_I\bigr\rangle=p!\binom{n}{p}pH.
\end{equation}
Integrating both (\ref{eq:sumdeltaomegaI2}) (after dividing by $\sigma_{n+1-p}$) and (\ref{eq:sumISpomegaI}) over $\partial M$, we deduce i\-ne\-qua\-li\-ty (\ref{eq:PoincareineqpformsboundaryinEuclideanspace}) from (\ref{eq:PoincareomegaI}).

If $M=\mathbb{D}^{n+1}$ is the $(n+1)$-dimensional flat disk standardly embedded in some $(n+1)$-dimensional affine subspace of $\mathbb{R}^{n+m}$, then $|H_0|=1=H$, $\sigma_{n+1-p}=n+1-p$ and  ${\rm Scal}^{\partial M}=n(n-1)$  along $\partial M=\mathbb{S}^n$ (the $n$-dimensional round sphere of sectional curvature $1$), so that (\ref{eq:PoincareineqpformsboundaryinEuclideanspace})  is an equality.

Conversely, if (\ref{eq:PoincareineqpformsboundaryinEuclideanspace}) is an equality, then for any tuple $I$, (\ref{eq:PoincareomegaI}) must be an equality.
For any $I$, we denote by $\hat{\omega}_I$ the $p$-form on $M$ such that ${\rm d}\hat{\omega}_I=0=\delta\hat{\omega}_I$ on $M$ and $\iota^*\hat{\omega}_I=\omega_I$ along $\partial M$; the existence and uniqueness of $\hat{\omega}_I$ is guaranteed by \cite[Theorem 2]{DuffSpencer52} and the vanishing of  $H_{\rm rel}^p(M)$,  see proof of Theorem \ref{t:Poincareineqpformsgeneral} above. Then Theorem \ref{t:Poincareineqpformsgeneral} implies that, for any $I$, the $p$-form $\hat{\omega}_I$ is parallel on $M$, that $\delta^{\partial M}\omega_I=-\sigma_{n+1-p}\nu\lrcorner\,\hat{\omega}_I$ and that $S^{[p-1]}(\nu\lrcorner\,\hat{\omega}_I)=(nH-\sigma_{n+1-p})\nu\lrcorner\,\hat{\omega}_I$ hold along $\partial M$. By (\ref{eq:sumdeltaomegaI2}), this implies
\begin{gather}\label{eq:sumdeltaomegaI2bis}
\sigma_{n+1-p}\sum_I|\nu\lrcorner\,\hat{\omega}_I|^2=\sum_I\frac{\bigl|\delta^{\partial M}\omega_I\bigr|^2}{\sigma_{n+1-p}}
=p!\binom{n}{p}p\left(\frac{n|H_0|^2-\frac{p-1}{n(n-1)}{\rm Scal}^{\partial M}}{\sigma_{n+1-p}}\right).
\end{gather}
Next we show that $\sum_I|\nu\lrcorner\,\hat{\omega}_I|^2$ is constant along $\partial M$. Differentiating along any $X\in T\partial M$ and using the fact that $\hat{\omega}_I$ is parallel on $M$, we have
\begin{gather*}
X\biggl(\frac{1}{2}\sum_I|\nu\lrcorner\,\hat{\omega}_I|^2\biggr)
=\sum_I\bigl\langle\nabla_X(\nu\lrcorner\,\hat{\omega}_I),\nu\lrcorner\,\hat{\omega}_I\bigr\rangle
=-\sum_I\langle SX\lrcorner\,\hat{\omega}_I,\nu\lrcorner\,\hat{\omega}_I\rangle\\
\qquad{}=-\sum_I\bigl\langle \iota^*(SX\lrcorner\,\hat{\omega}_I),\nu\lrcorner\,\hat{\omega}_I\bigr\rangle
=-\sum_I\langle SX\lrcorner\,\omega_I,\nu\lrcorner\,\hat{\omega}_I\rangle
=-\sum_I\bigl\langle (SX\lrcorner\,{\rm d}x_I)^T,\nu\lrcorner\,\hat{\omega}_I\bigr\rangle.
\end{gather*}

In the following, we will prove that the last sum vanishes.
For this, we will express the $(p-1)$-form $\nu\lrcorner\,\hat{\omega}_I$  in terms of the data of the immersion $\iota_0\colon\partial M\to\mathbb{R}^{n+m}$. That identity will be used at several places in the proof. We denote by $\mathbb{I}$ the second fundamental form of $\iota_0$. Using \cite[equation~(3.3)]{Savo05}, we have that, for each $p$-tuple $I$,
\begin{align}
\delta^{\partial M}\omega_I={} &\sum_{k=1}^p(-1)^{k+1}\mathbb{I}_{{\rm d}x_{i_k}^\perp}^{[p-1]}\bigl({\rm d}x_{i_1}^T\wedge \dots \wedge\widehat{{\rm d}x_{i_k}^T}\wedge\dots \wedge {\rm d}x_{i_p}^T\bigr)\nonumber\\
&{}-n \sum_{k=1}^p(-1)^{k+1}\bigl\langle H_0, {\rm d}x_{i_k}^\perp\bigr\rangle\,{\rm d}x_{i_1}^T\wedge\dots\wedge\widehat{{\rm d}x_{i_k}^T}\wedge\dots \wedge {\rm d}x_{i_p}^T\nonumber\\
 ={} &\sum_{k=1}^p\sum_{a=1}^m(-1)^{k+1}\langle {\rm d}x_{i_k},\nu_a\rangle\mathbb{I}_{\nu_a}^{[p-1]}\bigl({\rm d}x_{i_1}^T\wedge \dots \wedge\widehat{{\rm d}x_{i_k}^T}\wedge\dots \wedge {\rm d}x_{i_p}^T\bigr)\nonumber\\
 &{}-n \sum_{k=1}^p(-1)^{k+1}\langle H_0, {\rm d}x_{i_k}\rangle\,{\rm d}x_{i_1}^T\wedge\dots\wedge\widehat{{\rm d}x_{i_k}^T}\wedge\dots \wedge {\rm d}x_{i_p}^T\nonumber\\
 ={}&\sum_{a=1}^m\mathbb{I}_{\nu_a}^{[p-1]}\bigl((\nu_a\lrcorner\,{\rm d}x_I)^T\bigr)-(nH_0\lrcorner\,{\rm d}x_I)^T,\label{eq:deltaomegaISavo}
\end{align}
where $(\nu_a)_{1\leq a\leq m}$ is a pointwise o.n.b.\ of $T^\perp \partial M$ seen as a subspace of $\mathbb{R}^{n+m}$. Because of $\delta^{\partial M}\omega_I=-\sigma_{n+1-p}\nu\lrcorner\,\hat{\omega}_I$, we obtain
\begin{equation}\label{eq:nuinterioromega}
\nu\lrcorner\,\hat{\omega}_I=-\frac{1}{\sigma_{n+1-p}}\cdot\biggl(\sum_{a=1}^m\mathbb{I}_{\nu_a}^{[p-1]}\bigl((\nu_a\lrcorner\,{\rm d}x_I)^T\bigr)-(nH_0\lrcorner\,{\rm d}x_I)^T\biggr).
\end{equation}
Therefore, in order to show that $\sum_I\bigl\langle (SX\lrcorner\,{\rm d}x_I)^T,\nu\lrcorner\,\hat{\omega}_I\bigr\rangle=0$, it is sufficient to show that both sums $\sum_I\bigl\langle (SX\lrcorner\,{\rm d}x_I)^T,\sum_{a=1}^m\mathbb{I}_{\nu_a}^{[p-1]}\bigl((\nu_a\lrcorner\,{\rm d}x_I)^T\bigr)\bigr\rangle$ and $\sum_I\bigl\langle (SX\lrcorner\,{\rm d}x_I)^T,(nH_0\lrcorner\,{\rm d}x_I)^T\bigr\rangle$ vanish.
Let us make the computation for the first sum, the second can be done in the same way.
We denote by ${\rm e}_J={\rm e}_{j_1}\wedge\dots \wedge {\rm e}_{j_{p-1}}$ with $j_1<j_2<\dots<j_{p-1}$ the orthonormal frame of $\Lambda^{p-1}T_x^*\partial M$ induced by a local orthonormal frame $\{{\rm e}_1,\dots,{\rm e}_n\}$ of $T\partial M$.
For any $a=1,\dots,m$,
\begin{gather*}
\sum_I\bigl\langle (SX\lrcorner\,{\rm d}x_I)^T,\mathbb{I}_{\nu_a}^{[p-1]}\bigl((\nu_a\lrcorner\,{\rm d}x_I)^T\bigr)\bigr\rangle=\sum_{I,J}\bigl\langle (SX\lrcorner\,{\rm d}x_I)^T,{\rm e}_J\bigr\rangle \bigl\langle \mathbb{I}_{\nu_a}^{[p-1]}\bigl((\nu_a\lrcorner\,{\rm d}x_I)^T\bigr),{\rm e}_J\bigr\rangle\\
\hphantom{}=\sum_{I,J}\bigl\langle SX\lrcorner\,{\rm d}x_I,{\rm e}_J\bigr\rangle \bigl\langle \nu_a\lrcorner\,{\rm d}x_I,\mathbb{I}_{\nu_a}^{[p-1]}({\rm e}_J)\bigr\rangle
=\sum_{I,J}\bigl\langle {\rm d}x_I,SX\wedge {\rm e}_J\bigr\rangle \bigl\langle {\rm d}x_I,\nu_a\wedge \mathbb{I}_{\nu_a}^{[p-1]}({\rm e}_J)\bigr\rangle\\
\hphantom{}=\sum_{J}\bigl\langle SX\wedge {\rm e}_J,\nu_a\wedge \mathbb{I}_{\nu_a}^{[p-1]}({\rm e}_J)\bigr\rangle
=0,
\end{gather*}
because of $\nu_a\perp T\partial M$.
Similarly, $\sum_I\bigl\langle (SX\lrcorner\,{\rm d}x_I)^T,(nH_0\lrcorner\,{\rm d}x_I)^T\bigr\rangle=0$ by $H_0\perp T\partial M$.


If $\partial M$ is connected, which will be assumed from now on, then $\sum_I|\nu\lrcorner\,\hat{\omega}_I|^2$ is constant along~$\partial M$. By (\ref{eq:sumdeltaomegaI2bis}) and the assumption that (\ref{eq:PoincareineqpformsboundaryinEuclideanspace}) is an equality, this constant is given by
\begin{align*}
\sum_I|\nu\lrcorner\,\hat{\omega}_I|^2\int_{\partial M}\sigma_{n+1-p}\,{\rm d}\mu_g&= p!\binom{n}{p}p\int_{\partial M}\frac{n|H_0|^2-\frac{p-1}{n(n-1)}{\rm Scal}^{\partial M}}{\sigma_{n+1-p}}\,{\rm d}\mu_g\\
&= p!\binom{n}{p}p\int_{\partial M}H\,{\rm d}\mu_g,
\end{align*}
that is
\begin{equation}\label{eq:constantsumnuintomegaI2}
\sum_I|\nu\lrcorner\,\hat{\omega}_I|^2=p!\binom{n}{p}p\frac{\int_{\partial M}H\,{\rm d}\mu_g}{\int_{\partial M}\sigma_{n+1-p}\,{\rm d}\mu_g}.
\end{equation}
Injecting (\ref{eq:constantsumnuintomegaI2}) again  into (\ref{eq:sumdeltaomegaI2bis}), we deduce that
\begin{equation}\label{eq:pointwiseequalityPoincareineqpformsboundaryinEuclideanspace}
\frac{n|H_0|^2-\frac{p-1}{n(n-1)}{\rm Scal}^{\partial M}}{\sigma_{n+1-p}}=\frac{\int_{\partial M}H\,{\rm d}\mu_g}{\int_{\partial M}\sigma_{n+1-p}\,{\rm d}\mu_g}\cdot\sigma_{n+1-p}.
\end{equation}
Note that (\ref{eq:pointwiseequalityPoincareineqpformsboundaryinEuclideanspace}) holds in any codimension~$m$.

Next we look at the space of parallel forms on $M$. Let $\hat{\omega}:=\left(\hat{\omega}_I\right)_I\in\bigoplus_I\Omega^p(M)$. Fixing a~pointwise o.n.b.\ $({\rm e}_j)_{1\leq j\leq n}$ of $T\partial M$, the family
\[\left\{{\rm e}_J,\nu\wedge {\rm e}_K \mid 1\leq j_1<\dots< j_p\leq n,\;1\leq k_1<\dots<k_{p-1}\leq n \right\}\]
is a pointwise o.n.b.\ of $\Lambda^pT^*M$. Decomposing each $\hat{\omega}_I$ in that pointwise basis and the canonical basis $({\rm d}x_I)_I$ of $\Lambda^p(\mathbb{R}^{n+m})^*$ (where the $p$-tuples $I$ are ordered) respectively allows us to consider~$\hat{\omega}$ as a pointwise matrix with $\binom{n+m}{p}$ rows and $\binom{n+1}{p}$ columns.
Note that, because the pointwise linear map $\iota_0^*\colon\Lambda^p(\mathbb{R}^{n+m})^*\to\Lambda^pT^*\partial M$ is  surjective, the $\bigl(\omega_I=\iota_0^*({\rm d}x_I)\bigr)_I$ obviously span $\Lambda^pT^*\partial M$, which already shows that the $\binom{n}{p}$ first columns of the matrix $\hat{\omega}$, namely $\bigl(\hat{\omega}_I({\rm e}_J)\bigr)_{I,J}$, must be linearly independent since that matrix has $\binom{n}{p}$ linearly independent rows. Next we would like to show that the rank of the whole matrix $\hat{\omega}$ is maximal, i.e., equal to $\binom{n+1}{p}$.

This already allows for finding expressions for the inner products of columns of the matrix~$\hat{\omega}$.
Namely, fix ${\rm e}_J$ and $\nu\wedge {\rm e}_K$ as above, then using equation \eqref{eq:nuinterioromega}, we compute
\begin{gather*}
\sum_I\hat{\omega}_I({\rm e}_J)\hat{\omega}_I(\nu\wedge {\rm e}_K)=\sum_I\langle\hat{\omega}_I,{\rm e}_J\rangle\cdot\langle\nu\lrcorner\,\hat{\omega}_I,{\rm e}_K\rangle\\
\hphantom{\sum_I\hat{\omega}_I({\rm e}_J)\hat{\omega}_I(\nu\wedge {\rm e}_K)}=-\frac{1}{\sigma_{n+1-p}}\!\sum_I\bigl\langle{\rm d}x_I,{\rm e}_J\bigr\rangle\biggl\langle\sum_{a=1}^m\mathbb{I}_{\nu_a}^{[p-1]}\bigl((\nu_a\lrcorner\,{\rm d}x_I)^T\bigr)-(nH_0\lrcorner\,{\rm d}x_I)^T\!, {\rm e}_K\biggr\rangle\\
\hphantom{\sum_I\hat{\omega}_I({\rm e}_J)\hat{\omega}_I(\nu\wedge {\rm e}_K)}=-\frac{1}{\sigma_{n+1-p}}\sum_I \bigl\langle {\rm d}x_I,{\rm e}_J\bigr\rangle\biggl\langle {\rm d}x_I,\sum_{a=1}^m\nu_a\wedge\mathbb{I}_{\nu_a}^{[p-1]}{\rm e}_K-nH_0\wedge {\rm e}_K\biggr\rangle\\
\hphantom{\sum_I\hat{\omega}_I({\rm e}_J)\hat{\omega}_I(\nu\wedge {\rm e}_K)}=-\frac{1}{\sigma_{n+1-p}}\biggl\langle {\rm e}_J,\sum_{a=1}^m \nu_a\wedge\mathbb{I}_{\nu_a}^{[p-1]}{\rm e}_K-nH_0\wedge {\rm e}_K \biggr\rangle
=0
\end{gather*}
because of both $\nu_a,H_0\perp T\partial M$. This shows that every among the $\binom{n}{p-1}$ last columns of $\hat{\omega}$, corresponding to the matrix $\bigl(\hat{\omega}_I(\nu\wedge {\rm e}_K)\bigr)_{I,K}$, is pointwise orthogonal to any of the $\binom{n}{p}$ first ones which correspond to the full-ranked matrix $\bigl(\hat{\omega}_I({\rm e}_J)\bigr)_{I,J}$.

We now look at the rank of the matrix $\bigl(\hat{\omega}_I(\nu\wedge {\rm e}_K)\bigr)_{I,K}$. For any $(p-1)$-tuples $J$, $K$, we compute
\begin{gather*}
\sum_I\hat{\omega}_I(\nu\wedge {\rm e}_J)\hat{\omega}_I(\nu\wedge {\rm e}_K)=\sum_I\langle\nu\lrcorner\,\hat{\omega}_I,{\rm e}_J\rangle\langle\nu\lrcorner\,\hat{\omega}_I,{\rm e}_K\rangle\\
 =\frac{1}{\sigma_{n+1-p}^2}\sum_I \biggl\langle {\rm d}x_I,\sum_{a=1}^m \nu_a\wedge\mathbb{I}_{\nu_a}^{[p-1]}{\rm e}_J-nH_0\wedge {\rm e}_J\biggr\rangle\cdot \biggl\langle {\rm d}x_I,\sum_{a=1}^m \nu_a\wedge\mathbb{I}_{\nu_a}^{[p-1]}{\rm e}_K-nH_0\wedge {\rm e}_K\biggr\rangle\\
 =\frac{1}{\sigma_{n+1-p}^2}\biggl\langle\sum_{a=1}^m \nu_a\wedge\mathbb{I}_{\nu_a}^{[p-1]}{\rm e}_J-nH_0\wedge {\rm e}_J,
\sum_{a=1}^m \nu_a\wedge\mathbb{I}_{\nu_a}^{[p-1]}{\rm e}_K-nH_0\wedge {\rm e}_K\biggr\rangle\\
 =\frac{1}{\sigma_{n+1-p}^2}\biggl(\sum_{a=1}^m\bigl\langle \mathbb{I}_{\nu_a}^{[p-1]}{\rm e}_J,\mathbb{I}_{\nu_a}^{[p-1]}{\rm e}_K\bigr\rangle
-2\bigl\langle\mathbb{I}_{nH_0}^{[p-1]}({\rm e}_J),{\rm e}_K\bigr\rangle+n^2|H_0|^2\bigl\langle {\rm e}_J,{\rm e}_K\bigr\rangle\biggr).
\end{gather*}
Here we notice that, in the particular case, where all $\mathbb{I}_{\nu_a}$ are simultaneously diagonalizable, i.e., if $[\mathbb{I}_{\nu_a},\mathbb{I}_{\nu_b}]=0$ for all $a$, $b$, we can choose $({\rm e}_j)_{1\leq j\leq n}$ so as to simultaneously diagonalize all~$\mathbb{I}_{\nu_a}$.
Then it is easy to show that $\sum_I\hat{\omega}_I(\nu\wedge {\rm e}_J)\hat{\omega}_I(\nu\wedge {\rm e}_K)=0$ for all $J\neq K$.
However, the $\sum_I\hat{\omega}_I(\nu\wedge {\rm e}_J)^2$ cannot be shown to be positive.
This means that, even if $[\mathbb{I}_{\nu_a},\mathbb{I}_{\nu_b}]=0$ for all~$a$,~$b$, one column of the matrix $\bigl(\hat{\omega}_I(\nu\wedge {\rm e}_K)\bigr)_{I,K}$ may vanish, in which case $\hat{\omega}$ will not be of full rank.\looseness=-1

From now on we furthermore assume that $\iota_0(\partial M)\subset V$, where $V$ is an $(n+1)$-dimensional affine subspace of $\mathbb{R}^{n+m}$.
Choosing a pointwise o.n.b.\ $(\nu_a)_{1\leq a\leq m}$ of $T^\perp\partial M\subset\mathbb{R}^{n+m}$ such that $\nu_1\in V$ and $\nu_a\perp V$ for all $a\geq2$, we obviously have $\mathbb{I}_{\nu_a}=0$ for all $a\geq2$ since $V\subset\mathbb{R}^{n+m}$ is totally geodesic. Therefore, choosing $({\rm e}_j)_{1\leq j\leq n}$ as an eigenbasis for the endomorphism $\mathbb{I}_{\nu_1}$ of~$T\partial M$, we have $\mathbb{I}_{\nu_1}{\rm e}_j=\kappa_j {\rm e}_j$ for all $1\leq j\leq n$ and the sum $\sum_I\hat{\omega}_I(\nu\wedge {\rm e}_J)\hat{\omega}_I(\nu\wedge {\rm e}_K)$ computed above simplifies to
\begin{gather*}
\sum_I\hat{\omega}_I(\nu\wedge {\rm e}_J)\hat{\omega}_I(\nu\wedge {\rm e}_K)\\
\qquad{}=\frac{1}{\sigma_{n+1-p}^2}\Big(\bigl\langle \mathbb{I}_{\nu_1}^{[p-1]}{\rm e}_J,\mathbb{I}_{\nu_1}^{[p-1]}{\rm e}_K\bigr\rangle
-2\langle\mathbb{I}_{nH_0}({\rm e}_J),{\rm e}_K\rangle+n^2|H_0|^2\langle {\rm e}_J,{\rm e}_K\rangle\Big)\\
\qquad{}=\frac{1}{\sigma_{n+1-p}^2}\biggl(\biggl(\sum_{j\in J}\kappa_j\biggr)^2-2\langle nH_0,\nu_1\rangle\biggl(\sum_{j\in J}\kappa_j\biggr)+\langle nH_0,\nu_1\rangle^2\biggr)\delta_{JK}\\
\qquad{}=\frac{1}{\sigma_{n+1-p}^2}\biggl(\sum_{j\in J}\kappa_j-\langle nH_0,\nu_1\rangle\biggr)^2\delta_{JK}
=\frac{1}{\sigma_{n+1-p}^2}\biggl(\sum_{j\notin J}\kappa_j\biggr)^2\delta_{JK},
\end{gather*}
where $\delta_{JK}=0$ if $J\neq K$ and $1$ if $J=K$. Now since $\iota_0\colon \partial M\to V$ is an isometric immersion of~an~$n$-dimensional manifold into an $(n+1)$-dimensional Euclidean space, there exists a point $x\in\partial M$ for which $\kappa_j(x)>0$ for all $1\leq j\leq n$ holds, see, e.g., \cite[p.~255]{GallotHulinLafontaine04}. At that point $x$, we can conclude that $\sum_I\hat{\omega}_I(\nu\wedge {\rm e}_J)\hat{\omega}_I(\nu\wedge {\rm e}_K)>0$ if $J=K$ and vanishes otherwise.
Therefore the columns of \smash{$\bigl(\hat{\omega}_I(\nu\wedge {\rm e}_K)\bigr)_{I,K}$ } form an orthogonal system of nonzero vectors at $x$, from which can be deduced that the whole matrix $\hat{\omega}$ has full rank $\binom{n+1}{p}$ at $x$. In turn, this implies that, at $x$, there are $\binom{n+1}{p}$ linearly independent rows in $\hat{\omega}$, that is $\binom{n+1}{p}$ linearly independent $\hat{\omega}_I$. Necessarily there must exist $\binom{n+1}{p}$ linearly independent parallel forms among the $\hat{\omega}_I$ on $M$, which is the maximal number allowed. As a first consequence, $\bigl(M^{n+1},g\bigr)$ must be flat (remember that $1\leq p\leq n$).
As a second consequence, at each point of $\partial M$, the family $(\nu\lrcorner\,\hat{\omega}_I)_I$ must span $\Lambda^{p-1}T^*\partial M$, so that $\iota\colon\partial M\to M$ must be totally umbilical by the identity $S^{[p-1]}(\nu\lrcorner\,\hat{\omega}_I)=(nH-\sigma_{n+1-p})\nu\lrcorner\,\hat{\omega}_I$ for all $I$ and the assumption $p\geq2$, see Remark \ref{r:LambdapTMtrivial}.
Since $M$ is flat and $\iota$ is totally umbilical in the Einstein manifold $M$, the mean curvature $H$ must be constant --  and positive because of $(n+1-p)H\geq\sigma_{n+1-p}>0$. Up to rescaling $g$ on $M$ as well as the Euclidean metric on $\mathbb{R}^{n+m}$, it may be assumed that $H=1$ along $\partial M$. By \cite[Theorem 13]{RaulotSavo11}, because $\bigl(M^{n+1},g\bigr)$ is flat, $\iota$ is totally umbilical and with constant mean curvature $1$, the manifold $\bigl(M^{n+1},g\bigr)$ must be isometric to the $(n+1)$-dimensional flat disk.
Moreover, identity (\ref{eq:pointwiseequalityPoincareineqpformsboundaryinEuclideanspace}) implies that $|H_0|=1=H$ along~$\partial M$. But then  by the proof of the equality case in \cite[Theorem 1.2]{MiaoWang14}, there exists  an isometric immersion $M\to V$ extending $\iota$. Since $\partial M\cong \mathbb{S}^n$ has constant sectional curvature $1$, the immersion $\iota_0$ must be an embedding by standard results due to Hadamard and Cohn--Vossen, see, e.g., \cite{doCarmoWarner70}. Again  by the proof of \cite[Theorem~1.2]{MiaoWang14},  the above isometric immersion $M\to V$ extending $\iota$  is  an~embedding. This shows that $\bigl(M^{n+1},g\bigr)$ is isometric to the $(n+1)$-dimensional flat disk standardly embedded in $V$. This concludes the proof of Theorem \ref{t:PoincareineqpformsboundaryinEuclideanspace}.
\end{proof}
\begin{thebibliography}{99}
\bibitem[3]{doCarmoWarner70}
do~Carmo M.P., Warner F.W., Rigidity and convexity of hypersurfaces in spheres,
  \href{https://doi.org/10.4310/jdg/1214429378}{\textit{J.~Differential Geometry}} \textbf{4} (1970), 133--144.

\bibitem[4]{DuffSpencer52}
Duff G.F.D., Spencer D.C., Harmonic tensors on {R}iemannian manifolds with
  boundary, \href{https://doi.org/10.2307/1969771}{\textit{Ann. of Math.}} \textbf{56} (1952), 128--156.

\bibitem[6]{GallotHulinLafontaine04}
Gallot S., Hulin D., Lafontaine J., Riemannian geometry, 3rd~ed., \textit{Universitext},
  \href{https://doi.org/10.1007/978-3-642-18855-8}{Springer}, Berlin, 2004.

\bibitem[13]{MiaoWang14}
Miao P., Wang X., Boundary effect of {R}icci curvature, \href{http://doi.org/10.4310/jdg/1460463563}{\textit{J.~Differential
  Geom.}} \textbf{103} (2016), 59--82, \href{https://arxiv.org/abs/1408.2711}{arXiv:1408.2711}.

\bibitem[14]{RaulotSavo11}
Raulot S., Savo A., A {R}eilly formula and eigenvalue estimates for
  differential forms, \href{https://doi.org/10.1007/s12220-010-9161-0}{\textit{J.~Geom. Anal.}} \textbf{21} (2011), 620--640,
  \href{https://arxiv.org/abs/1003.0817}{arXiv:1003.0817}.

\bibitem[16]{Savo05}
Savo A., On the first {H}odge eigenvalue of isometric immersions, \href{https://doi.org/10.1090/S0002-9939-04-07702-0}{\textit{Proc.
  Amer. Math. Soc.}} \textbf{133} (2005), 587--594.

\end{thebibliography}
\end{document}
